% ============================================================
%  SEÇÃO 3 — METODOLOGIA
%
%  A metodologia descreve COMO você vai conduzir a pesquisa.
%  Deve ser clara o suficiente para alguém reproduzir seu trabalho.
%
%  Estrutura obrigatória (Art. 32º do Regulamento):
%    3.1 Tipo e natureza da pesquisa
%    3.2 Abordagem metodológica
%    3.3 Procedimentos e instrumentos
%    3.4 Ambiente e participantes (se houver)
%    3.5 Tecnologias e ferramentas (para projetos de software)
%
%  Tamanho sugerido: 3 a 6 páginas
% ============================================================

\chapter{METODOLOGIA}
\label{cap:metodologia}

% ----------------------------------------------------------
%  ESCOLHA SEU TIPO DE PESQUISA
%  Leia e mantenha apenas o que se aplica ao seu projeto.
%  Apague os comentários explicativos após escolher.
% ----------------------------------------------------------

% ----------------------------------------------------------
% 3.1 CLASSIFICAÇÃO DA PESQUISA
% ----------------------------------------------------------

\section{Classificação da Pesquisa}

% %->  Descreva como sua pesquisa se classifica.
%     Use o parágrafo-modelo abaixo e adapte conforme seu caso.

Esta pesquisa classifica-se, quanto à sua natureza, como [aplicada / básica]; quanto aos seus objetivos, como [exploratória / descritiva / explicativa]; quanto à abordagem do problema, como [quantitativa / qualitativa / misto]; e quanto aos procedimentos técnicos, como [estudo de caso / experimento / survey / pesquisa bibliográfica / DSR].

% ----------------------------------------------------------
%  TABELA DE REFERÊNCIA — apague após usar
% ----------------------------------------------------------
%
%  NATUREZA:
%    Aplicada    $\rightarrow$ solucionar um problema concreto
%    Básica      $\rightarrow$ gerar conhecimento sem aplicação imediata
%
%  OBJETIVOS:
%    Exploratória  $\rightarrow$ pouco se sabe sobre o tema; explorar
%    Descritiva    $\rightarrow$ descrever características de um fenômeno
%    Explicativa   $\rightarrow$ identificar causas e efeitos
%
%  ABORDAGEM:
%    Quantitativa $\rightarrow$ dados numéricos, métricas, estatísticas
%    Qualitativa  $\rightarrow$ entrevistas, análise de conteúdo, observação
%    Misto        $\rightarrow$ combina as duas
%
%  PROCEDIMENTO:
%    Estudo de Caso     $\rightarrow$ investigação aprofundada de 1 caso real
%    Experimento        $\rightarrow$ comparar soluções em condições controladas
%    Survey             $\rightarrow$ questionário com amostra de pessoas
%    Design Science     $\rightarrow$ criar e avaliar um artefato (sistema, modelo)
%    Pesquisa-ação      $\rightarrow$ pesquisador participa da solução
%    Bibliográfica      $\rightarrow$ análise de literatura existente


% ----------------------------------------------------------
% 3.2 PROCEDIMENTOS METODOLÓGICOS
% Descreva as etapas da sua pesquisa em ordem
% ----------------------------------------------------------

\section{Procedimentos Metodológicos}

% %->  Descreva as etapas que você vai seguir.
%     Liste em ordem cronológica.

A pesquisa será conduzida em [X] etapas, descritas a seguir:

\begin{enumerate}

  \item \textbf{[Nome da Etapa 1]:}
    [Descreva o que será feito nesta etapa, como será feito
    e qual o resultado esperado.]

  \item \textbf{[Nome da Etapa 2]:}
    [Descreva o que será feito nesta etapa.]

  \item \textbf{[Nome da Etapa 3]:}
    [Descreva o que será feito nesta etapa.]

  % Adicione mais etapas conforme necessário

\end{enumerate}

% ----------------------------------------------------------
%  MODELOS DE ETAPAS — copie e adapte conforme seu projeto
% ----------------------------------------------------------
%
%  Para projeto de SOFTWARE:
%  \item \textbf{Levantamento de Requisitos:}
%    Será realizado por meio de entrevistas semiestruturadas
%    com [X] usuários do sistema, com o objetivo de identificar
%    os requisitos funcionais e não funcionais.
%
%  \item \textbf{Modelagem e Prototipagem:}
%    Com base nos requisitos levantados, serão elaborados os
%    diagramas de caso de uso, diagrama de classes e
%    protótipo de interface no Figma.
%
%  \item \textbf{Implementação:}
%    O sistema será desenvolvido utilizando [tecnologias].
%    O código será versionado no GitHub.
%
%  \item \textbf{Avaliação de Usabilidade:}
%    Será aplicado o questionário SUS (System Usability Scale)
%    com [X] usuários para avaliar a usabilidade do sistema.
%
%  Para projeto de PESQUISA:
%  \item \textbf{Revisão Sistemática da Literatura:}
%    Será conduzida uma busca nas bases IEEE Xplore, ACM DL
%    e Periódicos CAPES com os termos [seus termos de busca].
%
%  \item \textbf{Coleta de Dados:}
%    [Descreva seu instrumento: questionário, experimento, etc.]
%
%  \item \textbf{Análise dos Dados:}
%    [Descreva como os dados serão analisados.]


% ----------------------------------------------------------
% 3.3 AMBIENTE E PARTICIPANTES
% (Pule esta seção se não tiver participantes humanos)
% ----------------------------------------------------------

\section{Ambiente e Participantes}

% %->  Descreva onde e com quem a pesquisa será realizada.
%     Se for pesquisa com pessoas, descreva o perfil dos participantes.
%     Se for experimento computacional, descreva o ambiente de testes.

[Descreva o local/contexto da pesquisa e o perfil dos participantes
ou o ambiente computacional utilizado.]

% Exemplo:
% A pesquisa será conduzida no contexto de [empresa/instituição],
% que possui [N] funcionários na área de [área]. Os participantes
% serão selecionados por conveniência, considerando os critérios
% de inclusão: [critério 1], [critério 2]. Estima-se uma amostra
% de [N] participantes.


% ----------------------------------------------------------
% 3.4 TECNOLOGIAS E FERRAMENTAS
% (Para projetos de desenvolvimento de software)
% ----------------------------------------------------------

\section{Tecnologias e Ferramentas}

% %->  Liste as tecnologias que você pretende utilizar.
%     Justifique cada escolha em 1-2 frases.

O Quadro~\ref{qua:tecnologias} apresenta as principais tecnologias e ferramentas previstas para o desenvolvimento do projeto.

\begin{table}[htb]
\caption{Tecnologias e ferramentas previstas}
\label{qua:tecnologias}
\begin{tabularx}{\textwidth}{lXl}
  \toprule
  \textbf{Tecnologia/Ferramenta} & \textbf{Função no Projeto} & \textbf{Categoria} \\
  \midrule
  {[Tecnologia 1]} & {[Para que será usada]} & {[Front-end / Back-end / BD / DevOps]} \\
  {[Tecnologia 2]} & {[Para que será usada]} & \\
  {[Tecnologia 3]} & {[Para que será usada]} & \\
  {[Tecnologia 4]} & {[Para que será usada]} & \\
  GitHub         & Versionamento do código-fonte & DevOps \\
  \bottomrule
\end{tabularx}
\fonte{Elaborado pelo autor (\oano).}
\end{table}

% Exemplos de tecnologias:
% React.js / Vue.js / Angular    $\rightarrow$ Front-end
% Node.js / Django / Spring Boot $\rightarrow$ Back-end
% PostgreSQL / MongoDB / MySQL   $\rightarrow$ Banco de Dados
% Docker / Kubernetes            $\rightarrow$ DevOps
% Figma                          $\rightarrow$ Prototipagem
% TensorFlow / scikit-learn      $\rightarrow$ Machine Learning


% ----------------------------------------------------------
% 3.5 CONSIDERAÇÕES ÉTICAS
% (Obrigatório se envolver pessoas — Art. 15º do Regulamento)
% ----------------------------------------------------------

% Descomente se sua pesquisa envolver participantes humanos:
%\section{Considerações Éticas}
%
% Esta pesquisa envolve seres humanos e, portanto, será submetida
% ao Comitê de Ética em Pesquisa (CEP) da instituição, conforme
% a Resolução CNS 466/2012. Os participantes serão informados
% sobre os objetivos da pesquisa e assinarão o Termo de
% Consentimento Livre e Esclarecido (TCLE) antes da coleta
% de dados. A identidade dos participantes será mantida em
% sigilo em todas as etapas da pesquisa.
